\documentclass[english,a4paper,12pt]{article}
%\documentclass[english,journal=jctcce,manuscript=article,articletitle=true,layout=twocolumn]{achemso}
\usepackage[english]{babel}
\usepackage[T1]{fontenc}
\usepackage[margin=2cm]{geometry}

\usepackage{doi}
\usepackage[version=4]{mhchem}
\usepackage{multicol}

\title{}
\author{}
\date{}

\begin{document}
\begin{multicols}{2}

\section{Methods}

\doi{10.1103/PRXQuantum.2.020310} qubit-ADAPT-VQE

\doi{10.1021/acs.jctc.4c00070} describe the explicitly correlated TC method and explore the use of optimized basis sets.

\doi{10.1021/acs.jpca.3c05882} incorporate an SCF step to ADAPT-VQE to optimize the orbitals.
The SCF step can be performed without considerable increase in the number of two-qubit gates.

\doi{0.1021/acs.jpclett.3c02536} utilize quantum information to optimize the active space and boost the convergence in CASSCF calculations.

\doi{10.48550/arXiv.2405.15422} present a measurement-cost effective method for inclusion of dynamical correlation effects.

\section{Experiments}

\doi{10.1038/nature23879} perform calculations with six qubits and hundreds of Pauli terms on molecules up to \ce{BeH2}.

\doi{10.1038/s41534-020-0259-3} perform VQE on \ce{H20} and without error mitigation they obtain accuracies comparable to chemical accuracy.

\doi{10.1126/science.abb9811} perform twelve qubit quantum chemical calculations and simulate the isomerization of diazene.
They variationally optimize the parametrized quantum circuit to prepare the HF state and perform error mitigation based on $N$-representability.

\doi{0.1021/acs.jctc.3c01281} they implement qubit-ADAPT-VQE in fiurst quantization.
They perform classical simulations to calculate \ce{H2} energy curves using five to seven qubits.
They use an explicitly corralated basis with up to 128 basis funcitons and reach chemical accuracy.

\section{Other}

\doi{10.1039/d1cs01184g} discuss basis sets and boundary conditions for electronic structure calculations on quantum computers.
They discuss that small, optimized numerical basis sets are appropriate.

\end{multicols}
\end{document}